\documentclass[12pt]{article}
\usepackage{amsmath,amssymb}

\begin{document}
\section*{{2019 AMC 10A Problems/Problem 21}}
Source: AoPS Wiki

\subsection*{Solution}
The triangle is placed on the sphere so that its three sides are tangent to the sphere. The cross-section of the sphere created by the plane of the triangle is also the incircle of the triangle. To find the inradius, use $\text{area} = \text{inradius} \cdot \text{semiperimeter}$ . The area of the triangle can be found by drawing an altitude from the vertex between sides with length $15$ to the midpoint of the side with length $24$ . The Pythagorean triple $9$ - $12$ - $15$ allows us easily to determine that the base is $24$ and the height is $9$ . The formula $\frac {\text{base} \cdot \text{height}} {2}$ can also be used to find the area of the triangle as $108$ , while the semiperimeter is simply $\frac {15 + 15 + 24} {2} = 27$ . After plugging into the equation, we thus get $108 = \text{inradius} \cdot 27$ , so the inradius is $4$ . Now, let the distance between $O$ and the triangle be $x$ . Choose a point on the incircle and denote it by $A$ . The distance $OA$ is $6$ , because it is just the radius of the sphere. The distance from point $A$ to the center of the incircle is $4$ , because it is the radius of the incircle. By using the Pythagorean Theorem, we thus find $x = \sqrt{6^2-4^2}=\sqrt{20} = \boxed{\textbf {(D) } 2 \sqrt {5}}$ .

\end{document}
